\chapter{Arquitetura do \emph{framework} e reprodutibilidade}\label{apendice-framework}

Este apêndice descreve como o código está organizado, o que cada execução registra e como um experimento pode ser reproduzido.
É material de apoio: o argumento e o método estão nos \hyperref[cap3-materiais-metodos]{Capítulos 3} a \hyperref[cap6-conclusoes]{6}.

\section{Organização do código}\label{apendice-organizacao}

O código está dividido em pacotes com responsabilidades separadas:

\begin{itemize}
  \tightlist
  \item
        \textbf{carregamento}: lê instâncias da TSPLIB e do banco de dados, valida o tipo de peso de aresta e rejeita coordenadas tridimensionais;
  \item
        \textbf{distâncias}: constrói e armazena as matrizes de distância, com arredondamento e conversão para inteiros de 32 bits segundo a convenção da TSPLIB;
  \item
        \textbf{algoritmos}: o algoritmo genético e as implementações de cada etapa --- inicialização, seleção, cruzamento, mutação, busca local e avaliação ---, com versões de CPU e de GPU atrás do mesmo contrato de chamada;
  \item
        \textbf{execução}: orquestra as repetições, grava os resultados parciais e retoma execuções interrompidas;
  \item
        \textbf{análise}: estatística descritiva, testes de hipótese, tabelas e diagramas de diferença crítica;
  \item
        \textbf{utilitários}: vetorização, medição de tempo, geração de números aleatórios e registro de mensagens.
\end{itemize}

A separação permite trocar a implementação de uma única etapa sem tocar nas demais, que é o que torna possível comparar as quatro variantes usando o mesmo algoritmo.

\section{Registros de execução}\label{apendice-registros}

Cada par algoritmo--instância produz um arquivo de resultados em formato textual, gravado ao final de \emph{cada} repetição, e não ao final do par.
A gravação é atômica: o conteúdo é escrito em um arquivo temporário, sincronizado em disco e só então renomeado sobre o destino.
Uma interrupção no meio de uma repetição deixa, portanto, o último estado consistente no lugar, e a execução pode ser retomada sem repetir trabalho já concluído e sem duplicar sementes.

O arquivo registra o resumo estatístico (número de execuções bem-sucedidas, média, desvio padrão, mínimo e máximo de tempo e de custo), as séries brutas por repetição (custo, tempo, \emph{gap}, semente e número de gerações executadas), a telemetria por fase descrita na \hyperref[apendice-telemetria]{Seção A.4} e os metadados da configuração: identificação da instância, soma de verificação do arquivo de origem, parâmetros do algoritmo e implementação usada em cada etapa.

\section{Base de dados de resultados}\label{apendice-banco}

Os registros são consolidados em um banco DuckDB, com uma tabela principal de execuções e tabelas filhas para as séries por repetição.
A consolidação é idempotente pela chave natural formada por instância e algoritmo: reimportar o mesmo resultado substitui o registro anterior dentro de uma única transação, em vez de acumular duplicatas.
Isso permite reimportar toda a pasta de resultados depois de uma recoleta sem limpar o banco manualmente, e garante que uma falha durante a escrita não deixe o banco em estado parcial.

\section{Telemetria por fase}\label{apendice-telemetria}

Cada repetição registra o tempo de oito etapas do ciclo evolutivo --- inicialização, seleção, cruzamento, mutação, avaliação, busca local, sobrevivência e sincronização com a CPU ---, além do volume transferido em cada direção entre CPU e GPU e do número de lançamentos de \emph{kernel}.
Os tempos são medidos no relógio da CPU, o que inclui o custo de lançamento das operações assíncronas e é a métrica relevante para comparar variantes do ponto de vista de quem executa o algoritmo; a medição por eventos no dispositivo existe como opção, mas impõe uma barreira de sincronização a cada chamada e por isso não é usada na coleta.
A soma das fases é comparada ao tempo total de parede de cada configuração, e a diferença residual mede o custo que não é atribuído a nenhuma etapa.

\section{Roteiro de reprodução}\label{apendice-reproducao}

A reprodução de um experimento segue cinco passos:

\begin{enumerate}
  \def\labelenumi{\arabic{enumi}.}
  \tightlist
  \item
        obter as instâncias da TSPLIB e carregá-las no banco de instâncias, que valida formato e tipo de peso de aresta na entrada;
  \item
        escolher a coorte, as variantes e o número de repetições;
  \item
        executar o orquestrador, que grava um arquivo por par algoritmo--instância e retoma do último ponto em caso de interrupção;
  \item
        consolidar os resultados no banco de resultados;
  \item
        gerar as tabelas e os diagramas a partir do banco consolidado.
\end{enumerate}

As sementes são determinísticas: cada repetição usa a semente base somada ao índice da repetição, de modo que a repetição \(k\) de um par pode ser reproduzida isoladamente.
O ambiente de execução --- versões de sistema operacional, biblioteca de GPU, compilador e placa --- é registrado junto com os resultados, porque os tempos absolutos dependem dele.
